\begin{explanationbox}
	\textbf{\textcolor{CExplanation}{一句话直觉}}

	在同一个圆（或等圆）中，弧的长度与它所对的圆周角大小是一一对应的，弧相等则角相等，反之亦然。

	\medskip
	\textbf{\textcolor{CExplanation}{核心拆解}}
	\begin{description}[style=nextline, leftmargin=2.8em, labelsep=0.5em, itemsep=0.45em]
		\item[\textbf{要点一（弧角对应）：}] 等弧必须对应相等的圆周角；等圆周角必须对应相等的弧。这是同圆或等圆中的等价关系。
		\item[\textbf{要点二（同圆前提）：}] 该结论只适用于同圆或等圆，如果两圆半径不同，弧相等并不能保证圆周角相等。
	\end{description}

	\medskip
	\textbf{\textcolor{CExplanation}{几何本质（圆心角媒介）}}

	弧相等 $\Rightarrow$ 圆心角相等 $\Rightarrow$ 圆周角相等（圆周角等于圆心角的一半）。反之亦然，圆周角相等 $\Rightarrow$ 圆心角相等 $\Rightarrow$ 弧相等。因此该结论本质是圆心角定理的推论。

	\medskip
	\textbf{\textcolor{CExplanation}{代数意义（符号等价）}}

	记弧长为 $l$，圆心角为 $\theta$，半径 $R$，则 $l = R\theta$（$\theta$ 用弧度制），但高中通常不引入弧长公式，直接使用弧的度数（圆心角度数）代替。本质上是比例关系：弧的度数等于圆心角的度数，圆周角是其一半。

	\medskip
	\textbf{\textcolor{CExplanation}{考点价值（证明工具）}}
	\begin{itemize}[leftmargin=*, itemsep=0.5em, topsep=0.4em]
		\item \textbf{考法一（证明角等）：} 题中给出弧相等，直接推出圆周角相等。
		\item \textbf{考法二（证明弧等）：} 题中给出圆周角相等，直接推出弧相等。
		\item \textbf{考法三（综合转换）：} 与其它定理结合，实现角与弧的多次转换。
	\end{itemize}

	\medskip
	\textbf{\textcolor{CExplanation}{顿悟点}}

	\textbf{圆周角相等不一定直接推出弧相等，必须满足“同圆或等圆”前提。}

	\medskip
	\textbf{\textcolor{CExplanation}{使用场景}}
	\begin{itemize}[leftmargin=*, itemsep=0.5em, topsep=0.4em]
		\item \textbf{场景一（圆内等角转换）：} 在圆中需要证明两个角相等时，若它们所对的弧相等，直接得证。
		\item \textbf{场景二（圆内等弧转换）：} 需要证明两段弧相等时，若它们所对的圆周角相等，直接得证。
		\item \textbf{场景三（多步推理）：} 在复杂图形中先通过其他步骤得到角或弧相等，再应用本结论继续推证。
	\end{itemize}
\end{explanationbox}
